\documentclass[aspectratio=169]{beamer}
\usepackage{listings}
\usepackage{xcolor}
\usepackage{booktabs}

\usetheme{Madrid}
\usecolortheme{seahorse}

\definecolor{codeblue}{rgb}{0,0,0.7}
\definecolor{backcolour}{rgb}{0.97,0.97,0.97}
\lstset{
  basicstyle=\ttfamily\small,
  backgroundcolor=\color{backcolour},
  keywordstyle=\color{codeblue},
  breaklines=true,
  numbers=none,
}

\title[Async Research Assistant]{Async Research Assistant}
\subtitle{AIENG Final Project — Topic 4}
\author{[Member 1] \and [Member 2] \and [Member 3]}
\institute{AI Engineering 110 · Spring 2026}
\date{May 23, 2026}

\begin{document}

% ---- Slide 1: Title --------------------------------------------------------
\begin{frame}
\titlepage
\end{frame}

% ---- Slide 2: Overview -----------------------------------------------------
\begin{frame}{Problem \& Solution}
\textbf{Problem}: Researching a question manually means opening multiple tabs,
reading several articles, and synthesising information yourself.

\vspace{1em}
\textbf{Our solution}: Ask once — get a cited answer from three authoritative sources.

\vspace{1em}
\begin{columns}
  \column{0.45\textwidth}
  \textbf{Input}
  \begin{itemize}
    \item A natural-language question
    \item Optional: source filter (wiki, arxiv, web)
  \end{itemize}

  \column{0.45\textwidth}
  \textbf{Output}
  \begin{itemize}
    \item 3–6 sentence synthesised answer
    \item Numbered references with URLs
    \item Which sources failed (if any)
  \end{itemize}
\end{columns}
\end{frame}

% ---- Slide 3: Architecture -------------------------------------------------
\begin{frame}{Architecture}
\begin{center}
\ttfamily\small
\begin{tabular}{l}
User / CLI\\
\quad $\downarrow$ ResearchRequest (Pydantic)\\
\textbf{core/researcher.py}\\
\quad $\leftarrow$ services/cache.py (TTL lookup)\\
\quad $\rightarrow$ concurrency/orchestrator.py\\
\qquad $\parallel$ fetch\_wikipedia \quad fetch\_arxiv \quad fetch\_web\\
\qquad $\downarrow$ asyncio.gather + Semaphore\\
\textbf{services/ai\_service.py} (retries + logging)\\
\quad $\downarrow$ ai/ module (DO NOT MODIFY)\\
\quad \quad synthesize() $\rightarrow$ AnswerWithCitations\\
\end{tabular}
\end{center}

\vspace{0.5em}
\textbf{Key principle}: every \texttt{ai.*} call goes through \texttt{ai\_service.py} — never directly.
\end{frame}

% ---- Slide 4: AI Module Integration ----------------------------------------
\begin{frame}[fragile]{AI Module Integration}
\textbf{Retry + timeout wrapper} — applied to every \texttt{ai.*} call:

\begin{lstlisting}[language=Python]
@_make_retry(max_attempts=4)
async def fetch_wikipedia(self, query, *, client=None):
    async with asyncio.timeout(settings.wikipedia_timeout):
        return await ai.fetch_wikipedia(query, client=client)
\end{lstlisting}

\vspace{0.5em}
\begin{itemize}
  \item \textbf{Retried}: \texttt{ConnectionError}, \texttt{TimeoutError}, \texttt{ProviderError}
  \item \textbf{Not retried}: \texttt{ValueError} (bad input — would fail again)
  \item Exponential backoff: 1s $\to$ 2s $\to$ 4s $\to$ 8s (max 30s)
  \item WARNING log before each retry sleep
\end{itemize}

\vspace{0.5em}
Provider (Anthropic / OpenAI / Gemini) selected via \texttt{LLM\_PROVIDER} env var.
\end{frame}

% ---- Slide 5: Concurrency --------------------------------------------------
\begin{frame}[fragile]{Concurrency: asyncio.gather + Semaphore}
\begin{lstlisting}[language=Python]
sem = asyncio.Semaphore(settings.max_concurrent)

async def _guarded(src_name):
    async with sem:
        return await _fetch_one(src_name, question)

results = await asyncio.gather(
    *[_guarded(s) for s in active_sources],
    return_exceptions=True,  # one failure != all fail
)
\end{lstlisting}

\begin{center}
\begin{tabular}{lrr}
\toprule
Mode & Time (s) & Speedup \\
\midrule
Sequential & 2.655 & 1.0$\times$ \\
Concurrent & 1.003 & \textbf{2.6$\times$} \\
\bottomrule
\end{tabular}
\end{center}
N=5 questions, simulated latency. Real network: typically 2–4$\times$ speedup.
\end{frame}

% ---- Slide 6: Robustness ---------------------------------------------------
\begin{frame}[fragile]{Robustness: Retries \& Graceful Degradation}
\textbf{Failure injection test}:

\begin{lstlisting}[language=Python]
# Monkeypatch arXiv to always fail
new_fns["arxiv"] = lambda q, **kw: (_ for _ in ()).throw(
    ConnectionError("arXiv is down"))
\end{lstlisting}

\textbf{Result}: system still returned a valid answer from Wikipedia + web, and
\texttt{ResearchResult.sources\_failed} correctly listed \texttt{"arxiv"}.

\vspace{1em}
\textbf{Input validation} via Pydantic:
\begin{itemize}
  \item Empty/whitespace question $\to$ rejected before any network call
  \item Question $>$ 500 chars $\to$ rejected
  \item Unknown source name $\to$ CLI exits with error message
\end{itemize}

\textbf{Logging}: structured \texttt{extra=\{...\}} dicts. No \texttt{print()} in \texttt{src/}.
\end{frame}

% ---- Slide 7: Testing ------------------------------------------------------
\begin{frame}{Testing}
\begin{columns}
  \column{0.5\textwidth}
  \textbf{66 tests, 81\% coverage}\\(minimum required: 60\%)

  \vspace{0.7em}
  \begin{itemize}
    \item \texttt{test\_ai\_smoke.py} — provided, must pass
    \item \texttt{test\_service.py} — retry, types
    \item \texttt{test\_cache.py} — hit/miss/expiry
    \item \texttt{test\_concurrency.py} — semaphore, degradation
    \item \texttt{test\_researcher.py} — full pipeline
    \item \texttt{test\_cli.py} — bad input, JSON output
    \item \texttt{test\_models.py} — validation
  \end{itemize}

  \column{0.45\textwidth}
  \textbf{All offline}\\(disconnect Wi-Fi — tests still pass)

  \vspace{0.7em}
  \textbf{autouse fixture} patches:\\
  \begin{itemize}
    \item \texttt{ai.fetch\_wikipedia} $\to$ canned Sources
    \item \texttt{ai.fetch\_arxiv} $\to$ canned Sources
    \item Web provider $\to$ \texttt{FakeWebSearch}
    \item LLM $\to$ \texttt{FakeLLM}
  \end{itemize}
\end{columns}
\end{frame}

% ---- Slide 8: Docker -------------------------------------------------------
\begin{frame}[fragile]{Docker}
\textbf{Multi-stage Dockerfile} (bonus +1 pt):

\begin{lstlisting}
# Builder: install deps into /opt/venv
FROM python:3.12-slim AS builder
RUN python -m venv /opt/venv && pip install -r requirements.txt

# Runtime: copy venv only — no gcc, no pip cache
FROM python:3.12-slim
COPY --from=builder /opt/venv /opt/venv
\end{lstlisting}

\vspace{0.5em}
\begin{lstlisting}[language=bash]
# Build
docker build --platform linux/amd64 -t researcher .

# Run offline demo (no keys needed)
docker run researcher

# With PostgreSQL
docker compose up --build
\end{lstlisting}

Non-root \texttt{appuser} runs the process. \texttt{.dockerignore} excludes \texttt{.env}, \texttt{.git}, test artefacts.
\end{frame}

% ---- Slide 9: Live Demo ----------------------------------------------------
\begin{frame}[fragile]{Demo}
\begin{lstlisting}[language=bash]
# Offline — no API keys
python demo_ai.py --offline --limit 2

# Real providers
python -m researcher ask "What is photosynthesis?"

python -m researcher ask "quantum computing" \
    --sources wiki,arxiv

python -m researcher ask "latest fusion research" \
    --no-cache --json-output

# Benchmark
python scripts/bench.py --n 5 --offline
\end{lstlisting}
\end{frame}

% ---- Slide 10: Lessons Learned ---------------------------------------------
\begin{frame}{Lessons Learned \& Q\&A}
\textbf{What worked well}:
\begin{itemize}
  \item \texttt{return\_exceptions=True} in \texttt{asyncio.gather} — one timeout
    never kills the whole response.
  \item Keeping \texttt{ai/} untouched behind \texttt{ai\_service.py} — trivial to mock.
  \item \texttt{autouse} fixture for no-network tests — zero test-by-test boilerplate.
\end{itemize}

\vspace{0.7em}
\textbf{What we'd do differently}:
\begin{itemize}
  \item Redis-backed distributed cache (in-memory doesn't survive restarts).
  \item Streaming LLM responses via Streamlit UI (bonus feature).
  \item \texttt{asyncio.to\_thread()} for \texttt{synthesize()} to avoid blocking the event loop.
\end{itemize}

\vspace{1em}
\begin{center}
  \Large \textbf{Questions?}
\end{center}
\end{frame}

\end{document}
